\section{Resultados Obtenidos}

\subsection{Optimización de Procesos}
Las mejoras operativas logradas tras la implementación de Odoo se centran en eliminar los procesos analógicos y unificar las operaciones de venta y almacén:

\begin{itemize}
    \item \textbf{Sincronización en Tiempo Real y Cero Doble Digitación:} Cada transacción realizada en el Punto de Venta (POS) en el mostrador de la calle Deán Valdivia descuenta de manera automática e inmediata el stock del almacén central, consolidando al sistema como la \textit{"Única Fuente de Verdad"} de las existencias. Esto elimina por completo la necesidad de que el vendedor registre la venta en papel y luego la digite en una hoja de Excel al final del día.
    \item \textbf{Consulta Digital Inmediata de Stock:} Frente a las consultas de clientes en mostrador o provincias, los vendedores consultan la disponibilidad directamente en la pantalla de Odoo, eliminando la antigua y lenta práctica de subir al almacén a realizar un conteo visual físico de estanterías.
    \item \textbf{Ventas de Provincias Estructuradas:} El despliegue de un Catálogo Digital visual permite enviar la oferta de productos en tiempo real y con stock garantizado a clientes de Puno, Cusco y Moquegua. La venta remota deja de ser una respuesta reactiva de fotos improvisadas vía WhatsApp, convirtiéndose en un flujo ágil soportado por el inventario real del sistema.
    \item \textbf{Gestión Proactiva de Abastecimiento:} Gracias al módulo de Tableros (Dashboards), la gerencia ya no decide las compras de mercadería por "intuición histórica" o conteos visuales. Ahora cuenta con un Reporte Semanal Automático de "Productos Estrella" y de lento movimiento, optimizando el uso de capital de trabajo y evitando compras innecesarias o desabastecimientos de stock crítico.
\end{itemize}

\subsection{Impacto de la Productividad}
Los resultados obtenidos al cierre de la implementación del plan se alinean con las metas proyectadas para junio de 2026. A continuación, se detalla la evaluación del impacto en los indicadores clave:

\begin{table}[H]
\centering
\renewcommand{\arraystretch}{1.3}
\setlength{\tabcolsep}{4pt}
\scriptsize
\begin{tabular}{|p{2.6cm}|p{2.8cm}|p{2.8cm}|p{3.2cm}|p{3.2cm}|}
\hline
\rowcolor{gray!20}
\textbf{KPI} & \textbf{Línea Base (Abr 2026)} & \textbf{Meta (Jun 2026)} & \textbf{Valor Final (Jul 2026)} & \textbf{Fuente} \\ \hline
\textbf{Exactitud de Inventario (ERI)} & 0\% (Registro empírico en cuadernos) & 75\% de coincidencia en el sistema. & \textbf{78\%} (superó la meta inicial por el piloto) & Reporte comparativo de cierre de stock del POS. \\ \hline
\textbf{Tiempo de Respuesta (Provincias)} & 45 - 60 Minutos (Conteo visual). & < 5 Minutos (Consulta rápida). & \textbf{< 2 Minutos} (Consulta digital instantánea). & Logs de atención y capturas de pantalla. \\ \hline
\textbf{Omnicanalidad} & 0 Canales oficiales. & 2 Canales oficiales. & \textbf{2 Canales activos} (Facebook e Instagram con catálogo). & Enlaces a redes sociales oficiales integradas. \\ \hline
\textbf{Crecimiento de Alcance} & < 0 consultas/mes (informal). & > 5 consultas/mes. & \textbf{> 12 consultas/mes} (Impulsado por pauta publicitaria). & Meta Business Suite y WhatsApp Business. \\ \hline
\textbf{Reportabilidad} & Nula / Empírica (Memoria visual). & Reporte Automático de productos estrella. & \textbf{Logrado} (Generación automática todos los lunes). & Tableros de control (dashboards). \\ \hline
\end{tabular}
\end{table}

\subsection{Adopción del Sistema}
La transición del factor humano involucró directamente al equipo operativo de 10 a 12 colaboradores, presentando los siguientes hitos:
\begin{itemize}
    \item \textbf{Nivel Inicial y Resistencia:} Existía una brecha en cuanto a proactividad técnica y una resistencia pasiva inicial debida a la falta de costumbre en el uso de herramientas específicas de almacén. El inventario "en la cabeza" de los vendedores antiguos generaba dependencia operativa.
    \item \textbf{Estrategia de Adopción:} Se designó al Encargado de Ventas y Almacén como líder responsable del escaneo obligatorio con código de barras. El Gerente General lideró activamente la transición, priorizando la carga de productos más fáciles y promoviendo una cultura de "aprender del error o del fracaso" para fomentar el crecimiento del personal.
    \item \textbf{Abandono Total del Papel:} Al término del proyecto, el uso de cuadernos de control manual se eliminó por completo. El personal integró la pistola lectora de códigos de barras y el POS como su herramienta instintiva diaria de trabajo en el mostrador.
\end{itemize}

\subsection{Plan de Capacitaciones}
Para garantizar que la adopción fuera exitosa, el plan de entrenamiento práctico sumó un total aproximado de 16 horas presenciales:

\begin{table}[H]
\centering
\renewcommand{\arraystretch}{1.2}
\begin{tabular}{|p{10cm}|c|}
\hline
\rowcolor{gray!20}
\textbf{Taller Impartido} & \textbf{Horas} \\ \hline
Carga Inicial e Inventariado (Conteo y codificación estandarizada) & 4 \\ \hline
Operación del POS, cobros digitales y Ventas en Mostrador & 4 \\ \hline
Control Logístico, recepción con pistola y Entrada de Proveedores & 4 \\ \hline
Catálogo Omnicanal, vinculación y flujos en WhatsApp Business & 2 \\ \hline
Reportabilidad y Dashboards (Exclusivo para Gerencia) & 2 \\ \hline
\rowcolor{gray!10}
\textbf{Total} & \textbf{16} \\ \hline
\end{tabular}
\end{table}
